% ============================================================
%  Section 6.1 – Software Stack
% ============================================================
\section{Software Stack: ROS~2 Architecture, MoveIt~2 Configuration,
         and Hardware Driver Integration}
\label{sec:ch6_software}

The software stack that drives the collaborative robotic cell is organised in four tiers:
the ROS~2 middleware and execution layer, the MoveIt~2 motion-planning layer, the hardware
abstraction and driver layer, and the digital-twin synchronisation layer.
Figure~\ref{fig:ch6_sw_stack} shows how these tiers relate to one another and to the physical
hardware.

% ------------------------------------------------------------------
%  Figure: four-tier software stack (TikZ)
% ------------------------------------------------------------------
\begin{figure}[H]
\centering
\begin{tikzpicture}[
    font=\small\sffamily,
    tier/.style={
        draw, thick, rounded corners=5pt,
        minimum width=13cm, minimum height=1.05cm,
        align=center, text width=12.5cm
    },
    arr/.style={->, >=stealth, thick, gray},
    lbl/.style={font=\footnotesize\sffamily\itshape, text=gray, align=left}
]

% Tier 4 – application / supervisory (top)
\node[tier, fill=blue!20] (t4) at (0, 4.5)
    {\textbf{Tier~4 – Supervisory / Application}\\
     \texttt{UnifiedAssemblySupervisor} \quad
     \texttt{HybridMTCController} \quad
     Digital Twin runners};

% Tier 3 – planning
\node[tier, fill=green!20] (t3) at (0, 3.0)
    {\textbf{Tier~3 – Motion Planning}\\
     MoveIt~2 \texttt{move\_group} \quad
     MoveIt Task Constructor (MTC) \quad
     OMPL / KDL / BioIK};

% Tier 2 – middleware
\node[tier, fill=orange!20] (t2) at (0, 1.5)
    {\textbf{Tier~2 – ROS~2 Middleware}\\
     Fast~DDS (rmw\_fastrtps\_cpp) \quad
     \texttt{ros2\_control} \quad
     \texttt{robot\_state\_publisher} \quad
     \texttt{tf2}};

% Tier 1 – hardware drivers (bottom)
\node[tier, fill=red!15] (t1) at (0, 0.0)
    {\textbf{Tier~1 – Hardware Drivers}\\
     \texttt{abb\_rws\_client} (IRC5/RWS) \quad
     AR4 serial driver (\texttt{/dev/ttyACM0}) \quad
     \texttt{realsense2\_camera}};

% Arrows
\foreach \a/\b in {t1/t2, t2/t3, t3/t4}
    \draw[arr] (\a.north) -- (\b.south);

% Side annotation
\node[lbl, right=0.6cm of t4] {Python async\\nodes};
\node[lbl, right=0.6cm of t3] {C++ nodes};
\node[lbl, right=0.6cm of t2] {C++ / Python};
\node[lbl, right=0.6cm of t1] {C++ / Python};

\end{tikzpicture}
\caption{Four-tier software stack of the collaborative robotic cell.
         Arrows indicate the direction of data and command flow from hardware up to the
         application layer.}
\label{fig:ch6_sw_stack}
\end{figure}

% ------------------------------------------------------------------
\subsection{ROS~2 Architecture}
\label{subsec:ch6_ros2}
% ------------------------------------------------------------------

The entire system runs under \textbf{ROS~2 Jazzy Jalisco} on Ubuntu~24.04.  All nodes are
launched from a single parameterised entry point, \texttt{master\_launch.py}, which accepts
arguments that select the operating mode (simulation-only, real-hardware, or digital-twin) and
conditionally starts the appropriate subset of nodes.

\textbf{Node topology.} Figure~\ref{fig:ch6_node_topology} maps every ROS~2 node, its callback-group assignment, and
the topics and actions that connect them.  The topology deliberately separates
time-critical safety paths (the \texttt{ZoneDetectionManager}) from the high-level planning
logic (the \texttt{UnifiedAssemblySupervisor}) so that a slow MoveIt planning call on the
supervisor cannot block a zone-breach interrupt.

% ------------------------------------------------------------------
%  Figure: ROS 2 node topology (TikZ)
% ------------------------------------------------------------------
\begin{figure}[H]
    \centering
    \resizebox{\textwidth}{!}{%
    \begin{tikzpicture}[
        font=\scriptsize,
        node distance=0.7cm and 1.1cm,
        ext/.style={draw, dashed, rounded corners=2pt, align=center, fill=gray!10,
                    minimum width=2.5cm, minimum height=0.85cm},
        sup/.style={draw, rounded corners=2pt, align=center, fill=blue!10,
                    minimum width=2.9cm, minimum height=0.85cm},
        task/.style={draw, rounded corners=2pt, align=center, fill=orange!12,
                     minimum width=2.5cm, minimum height=0.85cm},
        plan/.style={draw, rounded corners=2pt, align=center, fill=green!12,
                     minimum width=2.5cm, minimum height=0.85cm},
        safe/.style={draw, rounded corners=2pt, align=center, fill=red!12,
                     minimum width=2.5cm, minimum height=0.85cm},
        twin/.style={draw, rounded corners=2pt, align=center, fill=cyan!10,
                     minimum width=2.5cm, minimum height=0.85cm},
        arrow/.style={-Latex, thick},
        action/.style={-Latex, thick, dashed, color=blue!70!black},
        service/.style={-Latex, thick, dotted, color=green!55!black},
        safety/.style={-Latex, thick, color=red!75!black}
    ]
        % Column 1: external inputs
        \node[ext]                          (yolo)      {YOLO detection};
        \node[ext, below=of yolo]           (realsense) {RealSense camera};
        \node[ext, below=of realsense]      (fireb)     {Firestore bridge / GUI};

        % Column 2: supervisor
        \node[sup, right=1.5cm of realsense] (sup) {Unified Assembly\\Supervisor\\(Python, asyncio)};

        % Column 3: task servers
        \node[task]                         (ar4srv) at ($(sup.east) + (2.2,  1.5)$) {AR4 task server};
        \node[task]                         (abbsrv) at ($(sup.east) + (2.2, -1.5)$) {ABB task server};

        % Column 4: planning + safety
        \node[plan, right=of ar4srv]        (mtc)    {Hybrid MTC\\controller};
        \node[safe, right=of abbsrv]        (zone)   {Zone detection\\manager};

        % Column 5: digital twin
        \node[twin, right=of mtc]           (jstraj) {joint\_state\_to\\\_trajectory.py};
        \node[twin, right=of zone]          (simrun) {program\_sim\\\_runner.py};
        \node[twin] (divmon) at ($(jstraj)!0.5!(simrun)$) {divergence\_monitor.py};

        % Inputs into supervisor (converge into left side of sup)
        \draw[arrow] (yolo.east)      -| ([xshift=-0.25cm]sup.north);
        \draw[arrow] (realsense.east) -- (sup.west);
        \draw[arrow] (fireb.east)     -| ([xshift=-0.25cm]sup.south);

        % Supervisor to task servers (actions)
        \draw[action] (sup.east) -- (ar4srv.west);
        \draw[action] (sup.east) -- (abbsrv.west);

        % Supervisor to MTC (service): route ABOVE the task-server row
        % (mirror of the safety arrow's clearance — must clear ar4srv's top edge)
        \draw[service] ([xshift=0.25cm]sup.north) -- ++(0,2.4) -| (mtc.north);

        % Zone safety preempt: route BELOW the task-server row
        \draw[safety] (zone.south) -- ++(0,-1.0) -| ([xshift=0.25cm]sup.south);

        % Task servers to planning / safety
        \draw[action] (ar4srv) -- (mtc);
        \draw[action] (abbsrv) -- (zone);

        % Twin paths
        \draw[arrow] (mtc)  -- (jstraj);
        \draw[arrow] (zone) -- (simrun);
        \draw[arrow] (jstraj.south) -- (divmon.north);

        \begin{scope}[on background layer]
            \node[draw, dashed, rounded corners=4pt,
                  fit=(yolo)(realsense)(fireb), inner sep=0.24cm,
                  label=above:External inputs] {};
            \node[draw, dashed, rounded corners=4pt,
                  fit=(ar4srv)(abbsrv), inner sep=0.24cm,
                  label=above:Task servers] {};
            \node[draw, dashed, rounded corners=4pt,
                  fit=(mtc)(zone), inner sep=0.24cm,
                  label=above:Planning \& safety] {};
            \node[draw, dashed, rounded corners=4pt,
                  fit=(jstraj)(simrun)(divmon), inner sep=0.24cm,
                  label=above:Digital twin] {};
        \end{scope}
    \end{tikzpicture}%
    }
    \caption{High-level ROS 2 node topology.}
    \label{fig:ch6_node_topology}
\end{figure}

\textbf{Executor and callback-group strategy.} Table~\ref{tab:ch6_cbgroups} summarises the callback-group assignment that prevents
mutual blocking between the supervisor's planning work and the safety watchdog's interrupt path.

\begin{table}[H]
\centering
\caption{Callback-group assignments in the \texttt{UnifiedAssemblySupervisor} node.}
\label{tab:ch6_cbgroups}
\setlength{\tabcolsep}{6pt}
\begin{tabular}{lll}
\toprule
\textbf{Callback / Client}          & \textbf{Group type}           & \textbf{Rationale} \\
\midrule
State-machine timer                 & MutuallyExclusive             & Sequential task flow \\
AR4 \& ABB action clients           & Reentrant                     & True parallel dispatch \\
Perception \& gripper service calls & Reentrant                     & Non-blocking queries \\
Zone-status subscriber              & MutuallyExclusive (separate)  & Interrupt isolation \\
ExecuteTask action servers (C++)    & MutuallyExclusive             & One goal at a time per arm \\
Gripper SetBool services (C++)      & MutuallyExclusive             & Sequential gripper control \\
\bottomrule
\end{tabular}
\end{table}

\textbf{Launch file architecture.} \texttt{master\_launch.py} accepts Boolean launch arguments
(\texttt{use\_sim}, \texttt{use\_real\_abb}, \texttt{use\_real\_ar4}, \texttt{twin\_mode})
that select exactly which nodes are activated.  This single entry-point design means that the
same workspace can be launched in six distinct configurations — from pure Gazebo simulation
through full dual-hardware with digital twin — without modifying any source file.

% ------------------------------------------------------------------
\subsection{MoveIt~2 Configuration}
\label{subsec:ch6_moveit}
% ------------------------------------------------------------------

MoveIt~2 provides the motion-planning backbone for both arms.  The configuration was generated
by the MoveIt Setup Assistant and then hand-tuned to match the physical constraints of the
heterogeneous cell.

\textbf{Robot descriptions.} Both manipulators are described in URDF/Xacro.  The combined
\texttt{dual\_arms\_twin.xacro} includes the ABB~IRB~120 macro, the AR4~Mk3 macro, the
custom gripper macros, and a shared table link that provides the common world frame.  Unique
robot-name prefixes (\texttt{abb\_} and \texttt{ar4\_}) are injected into every macro to
prevent the link-name collision that initially produced duplicate
\texttt{gripper\_base\_link} geometries and spurious self-collision reports.

\textbf{Semantic description (SRDF).} The Semantic Robot Description Format file declares two planning groups:

\begin{itemize}
    \item \texttt{irb120\_arm} — six revolute joints of the ABB~IRB~120, with
          \texttt{tool0} as the end-effector link.
    \item \texttt{ar\_manipulator} — six revolute joints of the AR4~Mk3, with
          \texttt{ar4\_ee\_link} as the end-effector link.
\end{itemize}

Separate gripper groups (\texttt{abb\_gripper} and \texttt{ar4\_gripper}) are declared for
open/close trajectory control.  The SRDF also lists the allowed-collision pairs regenerated
after the prefix fix; without these entries MoveIt's collision checker would report
false-positive collisions between the table and each arm's base link.

\textbf{Kinematics configuration.} Table~\ref{tab:ch6_ik} lists the IK solver assigned to each planning group.
The ABB arm uses KDL because its well-conditioned Jacobian and high repeatability make
numerical IK fast and reliable.  The AR4 uses BioIK, a population-based solver, because the
arm's geometric structure occasionally leads KDL into local minima near wrist singularities.

\begin{table}[H]
\centering
\caption{IK solver configuration per planning group.}
\label{tab:ch6_ik}
\setlength{\tabcolsep}{6pt}
\begin{tabular}{llll}
\toprule
\textbf{Group}          & \textbf{Solver} & \textbf{Timeout (s)} & \textbf{Attempts} \\
\midrule
\texttt{irb120\_arm}    & KDL             & 0.005                & 3 \\
\texttt{ar\_manipulator}& BioIK           & 0.050                & 5 \\
\texttt{abb\_gripper}   & KDL             & 0.005                & 1 \\
\texttt{ar4\_gripper}   & KDL             & 0.005                & 1 \\
\bottomrule
\end{tabular}
\end{table}

\textbf{OMPL planning pipeline.} Both arm groups use \textbf{RRTConnect} as the primary planner within OMPL, with a 1.0\,s
planning timeout for standard pick-and-place moves.  MTC escape-route computations use a
tighter 2.0\,s budget; if no valid 12-DOF resolution path is found within that window the
system defaults to a complete emergency shutdown to preserve hardware integrity.

\textbf{Trajectory execution controllers.} \texttt{ros2\_control} loads four trajectory controllers at runtime (Table~\ref{tab:ch6_ctrls}).
In digital-twin mode the \texttt{/twin} namespace prevents the Gazebo controllers from
conflicting with the real-hardware controllers sharing the same ROS~2 graph.

\begin{table}[H]
\centering
\caption{\texttt{ros2\_control} trajectory controllers and their namespace assignment.}
\label{tab:ch6_ctrls}
\setlength{\tabcolsep}{5pt}
\begin{tabular}{lll}
\toprule
\textbf{Controller name}                    & \textbf{Namespace}  & \textbf{Targets} \\
\midrule
\texttt{irb120\_controller}                 & \texttt{/}          & ABB 6-DOF arm \\
\texttt{ar4\_trajectory\_controller}        & \texttt{/}          & AR4 6-DOF arm \\
\texttt{abb\_gripper\_controller}           & \texttt{/}          & ABB gripper \\
\texttt{ar4\_gripper\_controller}           & \texttt{/}          & AR4 gripper \\
\texttt{twin/irb120\_controller}            & \texttt{/twin}      & Gazebo ABB twin \\
\texttt{twin/ar4\_trajectory\_controller}   & \texttt{/twin}      & Gazebo AR4 twin \\
\texttt{twin/abb\_gripper\_controller}      & \texttt{/twin}      & Gazebo ABB gripper \\
\texttt{twin/ar4\_gripper\_controller}      & \texttt{/twin}      & Gazebo AR4 gripper \\
\bottomrule
\end{tabular}
\end{table}

Figure~\ref{fig:ch6_moveit_pipeline} illustrates the complete motion-planning request lifecycle
from the supervisor's goal dispatch through trajectory execution on the physical or simulated
controllers.

% ------------------------------------------------------------------
%  Figure: MoveIt 2 planning pipeline (TikZ)
% ------------------------------------------------------------------
\begin{figure}[H]
    \centering
    \resizebox{\textwidth}{!}{%
    \begin{tikzpicture}[
        font=\scriptsize,
        node distance=0.8cm and 1.1cm,
        block/.style={draw, rounded corners=2pt, align=center,
                      minimum width=2.5cm, minimum height=0.85cm},
        start/.style={block, fill=blue!10},
        proc/.style={block, fill=green!10},
        exec/.style={block, fill=orange!12},
        fail/.style={block, fill=red!12},
        dec/.style={draw, diamond, aspect=1.6, align=center, inner sep=2pt,
                    minimum width=2.5cm, minimum height=0.9cm, fill=yellow!18},
        arrow/.style={-Latex, thick},
        lbl/.style={font=\tiny, fill=white, inner sep=1.5pt}
    ]
        \node[start]                  (goal)  {Supervisor\\sends goal};
        \node[proc,  right=of goal]   (scene) {Planning scene\\update (TF2)};
        \node[proc,  right=of scene]  (ompl)  {OMPL planner\\RRTConnect};
        \node[dec,   right=of ompl]   (valid) {Collision\\free?};
        \node[exec,  right=of valid]  (traj)  {Trajectory\\parameterised};
        \node[exec,  below=of traj]   (exec)  {\texttt{FollowJoint}\\\texttt{Trajectory}};
        \node[fail,  below=of valid]  (failb) {Return\\FAILURE};

        \draw[arrow] (goal)  -- (scene);
        \draw[arrow] (scene) -- (ompl);
        \draw[arrow] (ompl)  -- (valid);
        \draw[arrow] (valid) -- node[lbl, above]{Yes} (traj);
        \draw[arrow] (traj)  -- (exec);
        \draw[arrow] (valid) -- node[lbl, right]{No}  (failb);
        \draw[arrow] (failb.west) -| node[lbl, above, pos=0.25]{retry} (ompl.south);

        \begin{scope}[on background layer]
            \node[draw, dashed, rounded corners=4pt,
                  fit=(scene)(ompl)(valid), inner sep=0.24cm,
                  label=above:Planning stage] {};
            \node[draw, dashed, rounded corners=4pt,
                  fit=(traj)(exec), inner sep=0.24cm,
                  label=above:Execution stage] {};
        \end{scope}
    \end{tikzpicture}%
    }
    \caption{MoveIt 2 motion-planning request lifecycle.}
    \label{fig:ch6_moveit_pipeline}
\end{figure}

% ------------------------------------------------------------------
\subsection{Hardware Driver Integration}
\label{subsec:ch6_drivers}
% ------------------------------------------------------------------

Three hardware drivers bridge the ROS~2 control stack to the physical devices.

\textbf{ABB IRC5 via Robot Web Services (RWS).} The \texttt{abb\_rws\_client} package communicates with the IRC5 controller over HTTP/WebSocket
at \texttt{192.168.125.1}.  On connection it authenticates with the controller, subscribes to
the joint-measurement resource, and republishes joint positions as \texttt{sensor\_msgs/JointState}
messages on \texttt{/joint\_states} at 50\,Hz.  Outgoing trajectory goals are accepted by the
ABB task server, converted to RAPID \texttt{MoveAbsJ} instructions via the RWS execution
interface, and monitored through a blocking WebSocket confirmation.  The RWS interface also
exposes the IRC5 emergency-stop status, which is mirrored as a ROS~2 diagnostic topic.

\textbf{AR4 Mk3 via Teensy serial driver.} The AR4 driver node opens \texttt{/dev/ttyACM0} at 115200\,baud.  The communication protocol
uses ASCII line-delimited commands; trajectory targets are sent as
\texttt{GOTO\_JOINTS~j1~j2~j3~j4~j5~j6} strings, where each value is an angle in degrees.
The Teensy firmware acknowledges each command and streams encoder-derived joint angles back at
100\,Hz.  The driver converts these to \texttt{/joint\_states} messages and exposes a
\texttt{SetBool} service for gripper control.  Cartesian path planning
(\texttt{computeCartesianPath}) enforces a minimum 95\,\% linear fraction during the final
approach and retreat phases, preventing unpredictable joint-space arcs near cluttered brick
stacks.

\textbf{Intel RealSense D435i.} The \texttt{realsense2\_camera} launch file activates the colour stream at 1280\,$\times$\,720
at 30\,fps and the aligned depth stream at the same resolution.  A time-synchronised subscriber
in \texttt{real2sim\_bridge.py} pairs each YOLO bounding-box detection with the corresponding
aligned depth pixel to back-project the brick's 3D position into the world frame, which is then
used both to update the MoveIt planning scene and to spawn the matching brick model in the
Gazebo twin.

Figure~\ref{fig:ch6_driver_integration} summarises the data paths from each hardware device
to the ROS~2 application layer.

% ------------------------------------------------------------------
%  Figure: hardware driver data paths (TikZ)
% ------------------------------------------------------------------
\begin{figure}[H]
    \centering
    \resizebox{\textwidth}{!}{%
    \begin{tikzpicture}[
        font=\scriptsize,
        node distance=0.7cm and 1.0cm,
        hw/.style={draw, thick, rounded corners=2pt, align=center, fill=red!12,
                   minimum width=2.7cm, minimum height=0.85cm},
        drv/.style={draw, rounded corners=2pt, align=center, fill=orange!12,
                    minimum width=2.9cm, minimum height=0.85cm},
        ros/.style={draw, rounded corners=2pt, align=center, fill=gray!10,
                    minimum width=2.7cm, minimum height=0.85cm},
        app/.style={draw, rounded corners=2pt, align=center, fill=green!12,
                    minimum width=2.8cm, minimum height=0.85cm},
        arrow/.style={-Latex, thick}
    ]
        % Row 1: ABB
        \node[hw]                       (irc5)   {ABB IRC5};
        \node[drv, right=of irc5]       (rws)    {\texttt{abb\_rws\_client}\\HTTP / WebSocket};
        \node[ros, right=of rws]        (jsirc)  {\texttt{/joint\_states}\\50\,Hz};

        % Row 2: AR4
        \node[hw,  below=of irc5]       (teensy) {Teensy 4.1\\(AR4 Mk3)};
        \node[drv, right=of teensy]     (ar4drv) {\texttt{ar4\_driver}\\serial ASCII};
        \node[ros, right=of ar4drv]     (jsar4)  {\texttt{/joint\_states}\\100\,Hz};

        % Row 3: RealSense
        \node[hw,  below=of teensy]     (rs)     {RealSense D435i};
        \node[drv, right=of rs]         (rsdrv)  {\texttt{realsense2\_camera}};
        \node[ros, right=of rsdrv]      (rsout)  {\texttt{/camera/color}\\+ aligned depth};

        % Application aligned to middle row (jsar4) — same height as the rest;
        % top and bottom topics converge orthogonally into its north/south edges.
        \node[app, right=1.2cm of jsar4] (app) {Application\\(Supervisor / MoveIt / Twin)};

        \foreach \a/\b in {irc5/rws, rws/jsirc, teensy/ar4drv, ar4drv/jsar4, rs/rsdrv, rsdrv/rsout}
            \draw[arrow] (\a) -- (\b);

        \draw[arrow] (jsirc.east) -| (app.north);
        \draw[arrow] (jsar4.east) -- (app.west);
        \draw[arrow] (rsout.east) -| (app.south);

        \begin{scope}[on background layer]
            \node[draw, dashed, rounded corners=4pt,
                  fit=(irc5)(teensy)(rs), inner sep=0.25cm,
                  label=above:Hardware] {};
            \node[draw, dashed, rounded corners=4pt,
                  fit=(rws)(ar4drv)(rsdrv), inner sep=0.25cm,
                  label=above:Drivers] {};
            \node[draw, dashed, rounded corners=4pt,
                  fit=(jsirc)(jsar4)(rsout), inner sep=0.25cm,
                  label=above:ROS 2 topics] {};
        \end{scope}
    \end{tikzpicture}%
    }
    \caption{Hardware driver data paths from physical devices to the ROS 2 application layer.}
    \label{fig:ch6_driver_integration}
\end{figure}

\textbf{Sim-clock isolation for MoveIt.} When the Gazebo twin is active it enables the ROS~2 simulation clock
(\texttt{use\_sim\_time:~true}).  The \texttt{getCurrentState()} call inside MoveIt evaluates
IK timeouts against this clock; with a paused or non-unit-rate sim clock the timeout fires
immediately, causing all IK solves to fail.  The fix is to supply a real-time
\texttt{rclcpp::Clock} instance explicitly to the \texttt{PlanningSceneMonitor} constructor,
decoupling the IK solver from the simulated time source while still allowing Gazebo to control
the physics step rate independently.